\documentclass[12pt]{article}
\usepackage[utf8]{inputenc}
\usepackage[T1]{fontenc}
\usepackage{hyperref}
\title{Tajuasohn Language Resources for K-12 Curriculum Development}
\author{Compiled by Claude Sonnet 5 (claude-sonnet-5) for Andr\'as Kornai}
\date{}
\begin{document}
\maketitle
\section{Introduction}
Tajuasohn (ISO 639-3 tja) is a Kru language spoken in Sinoe County, Liberia, by five local clans. Speaker counts in the sources differ: 9,600 (1991) and about 20,000 (2020). The searches found only brief descriptions.

\section{Primary Educational Resources}
None found.

\section{Linguistic and Reference Materials}
\begin{itemize}
\item Wikipedia on Tajuasohn: \url{https://en.wikipedia.org/wiki/Tajuasohn_language}
\item Numeral-system data from the Max Planck Institute (EVA), crediting D. Drevlow, 1994: \url{https://lingweb.eva.mpg.de/channumerals/Tajuasohn.htm}
\end{itemize}

\section{Religious and Community Resources}
Joshua Project profile of the Tajuasohn in Liberia: \url{https://joshuaproject.net/people_groups/15203/LI}
No scripture or community organisation was confirmed.

\section{Digital and Audio Resources}
None found.

\section{Curriculum Development Context}
Essentially no teaching or reference materials were found. Any curriculum would require new documentation work in Sinoe County, with the community.
\end{document}